% Ensure included PDFs (banners) with newer versions embed without warnings
\pdfminorversion=7
\documentclass[aspectratio=169]{beamer}

\makeatletter
% Make LaTeX find theme .sty files in Template
\def\input@path{{Template/}}
\makeatother
\usepackage{graphicx}
% Make graphics (including banner PDFs) resolvable without TEXINPUTS
\graphicspath{{Template/}{Template/banners/}{../Common_Images/}}

%================================================================%
% Theme and Package Setup
%================================================================%
\usetheme{ucl}
\setbeamercolor{banner}{bg=darkpurple}

% Navigation and Footer
\setbeamertemplate{navigation symbols}{\vspace{-2ex}}
\setbeamertemplate{footline}[author title date]
\setbeamertemplate{slide counter}[framenumber/totalframenumber]

\usepackage[utf8]{inputenc}
\usepackage[british]{babel}
\usepackage{amsmath, amssymb, amsthm}
\usepackage{tikz}
\usepackage{pgfplots}
\pgfplotsset{compat=1.17}
\usetikzlibrary{calc, positioning, arrows.meta, shapes.geometric, trees, backgrounds, shapes.misc, graphs, quotes, shadows, fit, decorations.pathreplacing} 
\usefonttheme{professionalfonts}
\usepackage{eulervm}
\usepackage{listings}
\usepackage{xcolor}
\usepackage{algorithm}
\usepackage{algpseudocode}
\usepackage[square,authoryear]{natbib}

% Define custom colours
\makeatletter
\@ifundefined{color@stone}{%
    \definecolor{stone}{gray}{0.95}%
}{}
\makeatother
\definecolor{darkgreen}{rgb}{0.0, 0.5, 0.0}
\definecolor{uclblue}{cmyk}{1,0.24,0,0.64}

\setbeamercovered{transparent}

% Define theoremblock
\makeatletter
\newenvironment<>{theoremblock}[1]{%
    \begin{block}#2{#1}%
}{\end{block}}
\makeatother

% Section divider slides
\AtBeginSection[]{
  \begin{frame}
    \frametitle{\textbf{\Large\insertsectionhead}}
    \begin{center}
        \huge\textbf{\insertsectionhead}
    \end{center}
  \end{frame}
}

% Operator definitions
\DeclareMathOperator*{\argmin}{arg\,min}
\DeclareMathOperator{\prox}{prox}
\DeclareMathOperator{\proj}{proj}
\DeclareMathOperator{\divg}{div}

\title{Advanced Deep Learning for Inverse Problems}
\subtitle{From Algorithm Unrolling to Generative Models}
\author[Martin Benning (University College London)]{Martin Benning}
\date[CCMI]{CCMI -- Deep Learning in Theory and Practice \\[1cm] Tuesday, 19 May 2026}
\institute[]{University College London}

\begin{document}

% --- Title Frame ---
\begin{frame}
  \titlepage
\end{frame}

% --- Overview Frame ---
\begin{frame}{Lecture Overview}
  \begin{block}{Objectives for Today}
    \begin{itemize}
                \item Understand \textbf{null spaces} and enforce data consistency with \textbf{null-space networks}.
                \item Formulate inverse problems with \textbf{variational methods} and classical regularisation.
                \item Review \textbf{optimisation algorithms}: ISTA and PDHG.
                \item Cast learning as \textbf{bilevel optimisation}.
                \item Compute hyperparameter sensitivities via \textbf{IFT} and \textbf{Huber smoothing}.
                \item Learn regularisers with \textbf{NETT} and recover guarantees via \textbf{ICNNs}.
                \item Introduce \textbf{algorithm unrolling} (LISTA, LPD).
    \end{itemize}
  \end{block}
\end{frame}

\begin{frame}{Lecture Overview}
  \tableofcontents
\end{frame}

%================================================================%
\section{Recap: Inverse Problems \& End-to-End Learning}
%================================================================%

\begin{frame}{Recap: The Inverse Problem}
    Recall from Monday that we were dealing with linear inverse problems of the form
    \begin{align*}
        y^\delta = A x^\dagger + \eta \, ,
    \end{align*}
    where
    \begin{itemize}
        \item $x^\dagger \in \mathbb{R}^n$ is the unknown ground-truth object (e.g., an image).
        \item $A \in \mathbb{R}^{m \times n}$ is the known forward operator (physics of the measurement).
        \item $\eta \in \mathbb{R}^m$ is random measurement noise.
        \item $y^\delta \in \mathbb{R}^m$ is the observed data.
        \item $y = Ax^\dagger$ is the ideal (noise-free) data.
    \end{itemize}
\end{frame}

\begin{frame}{Recap: The Pitfalls of End-to-End Learning}
    Training a standard neural network $f_\theta$ to directly map $y^\delta \to x^\dagger$ (end-to-end learning) suffers from severe limitations:
    
    \vspace{0.5em}
    \begin{enumerate}
        \item \textbf{Adversarial Instability:} Because the inverse operator $A^{-1}$ (or $A^\dagger$) is ill-conditioned, the network must be learned to have a massive Lipschitz constant, making it highly sensitive to tiny perturbations.
        \item \textbf{Implicit Bias:} The network hallucinates details based on training data statistics, which can lead to disastrous generalisation failures on rare or atypical features.
        \item \textbf{Lack of Data Consistency:} The network output $\hat{x} = f_\theta(y^\delta)$ is not guaranteed to satisfy the physical measurements, meaning $A\hat{x} \neq y = Ax^\dagger$.
    \end{enumerate}
\end{frame}

\begin{frame}{Recap: Classical vs. End-to-End Learning}
    Recall this summary from yesterday's lecture:
    \begin{table}
        \centering
        \begin{tabular}{l|c|c}
             & \textbf{Classical Methods} & \textbf{Deep Learning (End-to-End)} \\
             \hline
             \textbf{Data Consistency} & Guaranteed ($Ax \approx y$) & Not Guaranteed \\
             \textbf{Stability} & Mathematically bounded & Unstable (High Lipschitz needed) \\
             \textbf{Assumptions} & Explicit (Hand-crafted) & Opaque (Training bias) \\
             \textbf{Image Quality} & Limited (Oversimplified) & High (Realistic textures) \\
        \end{tabular}
    \end{table}
    
    \vspace{1em}
    \textbf{Today's Goal:} We want to bridge this gap. We will combine the guaranteed data consistency and mathematical stability of classical methods with the superior image quality and data-driven priors of deep learning.
\end{frame}

%================================================================%
\section{Null-Space Networks \& Data Consistency}
%================================================================%

\begin{frame}{The Null Space of the Forward Operator}
    One fundamental reason inverse problems are difficult (ill-posed) is that the operator $A \in \mathbb{R}^{m \times n}$ often has a non-trivial \textbf{null space} (or kernel).
    
    \begin{block}{Definition: Null Space}
        The null space $\mathcal{N}(A)$ is the set of all vectors in $\mathbb{R}^n$ that the operator $A$ maps to zero, i.e.,
        \begin{align*}
            \mathcal{N}(A) = \{ v \in \mathbb{R}^n \mid Av = 0 \} \, .
        \end{align*}
    \end{block}
    
    If $v \in \mathcal{N}(A)$, then $A(x + v) = Ax + Av = Ax$. 
    
    \vspace{1em}
    This means the forward measurement process is \emph{completely blind} to any features that reside in the null space!
\end{frame}

\begin{frame}{Orthogonal Decomposition}
    By a \textbf{Fundamental Theorem of Linear Algebra}, the domain $\mathbb{R}^n$ is the direct sum of the null space and the row space of $A$, i.e.,
    \begin{align*}
        \mathbb{R}^n = \mathcal{N}(A) \oplus \mathcal{R}(A^\top) \, .
    \end{align*}
    Thus, any image $x \in \mathbb{R}^n$ can be orthogonally decomposed into two unique components $x = x_{\text{meas}} + x_{\text{null}}$, where
    \begin{itemize}
        \item $x_{\text{meas}} \in \mathcal{R}(A^\top)$: The part of the image that can actually be measured.
        \item $x_{\text{null}} \in \mathcal{N}(A)$: The part completely lost during the measurement process.
    \end{itemize}\pause
    
    \begin{block}{Why are they orthogonal?}
        If $x_{\text{null}} \in \mathcal{N}(A)$ and $x_{\text{meas}} = A^\top u \in \mathcal{R}(A^\top)$, their inner product is exactly zero, i.e.,\vspace{-0.15cm}
        \begin{align*}
            \langle x_{\text{null}}, x_{\text{meas}} \rangle = \langle x_{\text{null}}, A^\top u \rangle = \langle A x_{\text{null}}, u \rangle = \langle 0, u \rangle = 0 \, .
        \end{align*}

        \vspace{-0.15cm}
    \end{block}
\end{frame}

\begin{frame}{The Role of the Pseudo-Inverse}
    Recall from Monday that the Moore--Penrose pseudo-inverse $A^\dagger$ computes the \textbf{minimum norm} solution to the least squares problem $\min_x \|Ax - y\|^2$.
    
    \vspace{0.5em}
    Because $x = x_{\text{meas}} + x_{\text{null}}$ is an orthogonal decomposition, Pythagoras' theorem gives
    \begin{align*}
        \|x\|^2 = \|x_{\text{meas}}\|^2 + \|x_{\text{null}}\|^2 \, .
    \end{align*}
    
    To minimise $\|x\|$, the pseudo-inverse naturally sets the null-space component exactly to zero. 
    
    \vspace{0.5em}
    Crucially, this means $A^\dagger y$ exactly recovers the measurable component (assuming zero noise) but does not attempt to guess what was lost, yielding
    \begin{align*}
        A^\dagger (A x^\dagger) = x_{\text{meas}} \, .
    \end{align*}
\end{frame}

\begin{frame}{The Role of Deep Learning}
    If $A^\dagger y$ perfectly recovers $x_{\text{meas}}$ in the noise-free setting, what is the role of deep learning?\pause
    
    \vspace{0.5em}
    \begin{alertblock}{Idealised vs. Realistic Settings}
        In an ideal, noise-free scenario, the neural network should not waste capacity trying to invert $A$. Its sole purpose is to realistically hallucinate the missing information $x_{\text{null}}$ using prior knowledge learned from training data.
        
        \vspace{0.5em}
        However, in realistic settings with noisy data $y^\delta$, the network must \emph{also} suppress noise and stabilise the inversion process.
    \end{alertblock}
\end{frame}

\begin{frame}{Null-Space Networks}
    A \textbf{Null-Space Network} isolates the data-driven model entirely within the null space of the operator $A$.
    
    \begin{theoremblock}{Null-Space Architecture}
        Let $x_0 = A^\dagger y$ be the pseudo-inverse solution. The network prediction takes the form\vspace{-0.15cm}
        \begin{align*}
            f_\theta(y^\delta) = x_0 + (I - A^\dagger A) \mathcal{H}_\theta(x_0)
        \end{align*}

        \vspace{-0.15cm}
        where $\mathcal{H}_\theta$ is any deep architecture (e.g., a UNet) and $P_{\mathcal{N}} = (I - A^\dagger A)$ is the orthogonal projection operator onto the null space $\mathcal{N}(A)$.
    \end{theoremblock}
        
        \begin{center}
        \resizebox{0.55\textwidth}{!}{
        \begin{tikzpicture}[
            box/.style={draw, thick, fill=red!20, minimum width=2.2cm, minimum height=1.2cm, align=center, rounded corners},
            net/.style={draw, thick, fill=uclblue!20, minimum width=2.5cm, minimum height=1.2cm, align=center, rounded corners},
            sum/.style={draw, circle, thick, minimum size=0.6cm, align=center},
            arr/.style={->, thick, >=latex}
        ]
            \node (input) at (-2.5, 0) {$y$};
            \node[box] (pinv) at (0, 0) {Pseudo-Inverse\\$A^\dagger$};
            \node[coordinate] (split) at (2.5, 0) {};
            \node[net] (unet) at (5, -2) {Deep Network\\$\mathcal{H}_\theta$};
            \node[box] (proj) at (9.5, -2) {Null-Space Projection\\$(I - A^\dagger A)$};
            \node[sum] (add) at (9.5, 0) {$+$};
            \node (output) at (11.5, 0) {$f_\theta(y)$};
            
            \draw[arr] (input) -- (pinv);
            \draw[-] (pinv) -- (split);
            \draw[arr] (split) -- node[above] {$x_0$} (add);
            \draw[arr] (split) |- (unet.west);
            \draw[arr] (unet) -- node[above] {$\mathcal{H}_\theta(x_0)$} (proj);
            \draw[arr] (proj) -- node[right] {$x_{\text{null}}$} (add);
            \draw[arr] (add) -- (output);
            
            \fill (split) circle (2pt);
        \end{tikzpicture}
        }
        \end{center}
\end{frame}

\begin{frame}{Proof of Strict Data Consistency}
    The network is mathematically guaranteed to be data consistent (in the absence of noise). We apply the forward operator $A$ to the network's output to observe
    
    \begin{align*}
        A f_\theta(y) &= A \left( A^\dagger y^\delta + (I - A^\dagger A) \mathcal{H}_\theta(x_0) \right) \, , \\
        &= A A^\dagger y^\delta + A(I - A^\dagger A) \mathcal{H}_\theta(x_0) \, , \\
        &= y + (A - A A^\dagger A) \mathcal{H}_\theta(x_0) \, .
    \end{align*}
    By the properties of the pseudo-inverse, $A A^\dagger A = A$. Therefore, we have
    \begin{align*}
        A f_\theta(y) = y + (A - A) \mathcal{H}_\theta(x_0) = y \, .
    \end{align*}
\end{frame}

\begin{frame}{Properties of Null-Space Networks}
    \begin{itemize}
        \item \textbf{Safety:} The network $\mathcal{H}_\theta$ can be arbitrarily expressive, deep, and non-linear. However, the projection $(I - A^\dagger A)$ ensures it \emph{physically cannot} alter the measured data.
        \item \textbf{Interpretability:} If the network hallucinates an artefact, we know with 100\% certainty that the artefact lies entirely in the null space of the measurement operator.
    \end{itemize}
    
    \vspace{0.5em}
    \textbf{Example: Inpainting.} If $A$ is a masking operator that hides the centre of an image, $A^\dagger A$ is the identity outside the mask and $0$ inside. $(I - A^\dagger A)$ zeroes out everything outside the mask. The network can only paint inside the missing region.
    
    \vspace{0.5em}
    \begin{center}
    \begin{tikzpicture}[scale=0.8, transform shape]
        % A^\dagger y
        \fill[uclblue!40] (0,0) rectangle (1.5,1.5);
        \fill[white] (0.5,0.5) rectangle (1.0,1.0);
        \draw[step=0.5cm, thick] (0,0) grid (1.5,1.5);
        \node at (0.75, 0.75) {\textbf{?}};
        \node[below] at (0.75, -0.2) {$A^\dagger y$};

        \node at (2.25, 0.75) {$+$};

        % (I - A^\dagger A) H
        \fill[white] (3,0) rectangle (4.5,1.5);
        \fill[red!40] (3.5,0.5) rectangle (4.0,1.0);
        \draw[step=0.5cm, thick] (3,0) grid (4.5,1.5);
        \node[below, align=center] at (3.75, -0.2) {$(I - A^\dagger A)\mathcal{H}_\theta(x_0)$};

        \node at (5.25, 0.75) {$=$};

        % f_theta(y)
        \fill[uclblue!40] (6,0) rectangle (7.5,1.5);
        \fill[red!40] (6.5,0.5) rectangle (7.0,1.0);
        \draw[step=0.5cm, thick] (6,0) grid (7.5,1.5);
        \node[below] at (6.75, -0.2) {$f_\theta(y)$};
    \end{tikzpicture}
    \end{center}
\end{frame}

\begin{frame}{Limitations of Null-Space Networks}
    While theoretically beautiful, Null-Space Networks have a major practical limitation:
    
    \begin{alertblock}{Computational Cost}
        Computing the pseudo-inverse $A^\dagger y$ and the projection $(I - A^\dagger A)$ requires computing the inverse of $A^\top A$. 
    \end{alertblock}
    
    For small problems, this is fine. For large-scale inverse problems in medical imaging (like 3D MRI or CT where $x \in \mathbb{R}^{256^3}$), explicitly constructing or inverting $A^\top A$ is computationally intractable.
\end{frame}

\begin{frame}{Regularised Null-Space Networks}
    To address computational intractability and handle noise $y^\delta$, we can replace the pseudo-inverse $A^\dagger$ with a stable, classical regularised reconstruction operator $R$.
    
    \vspace{0.5em}
    Replacing $A^\dagger$ with $R$ in our previous formulation yields
    \begin{align*}
        f_\theta(y^\delta) = R(y^\delta) + \mathcal{H}_\theta(R(y^\delta)) - R \left( A \mathcal{H}_\theta(R(y^\delta)) \right) \, .
    \end{align*}
    
    \begin{itemize}
        \item $R(y^\delta)$ provides a stable, noise-suppressed initial reconstruction.
        \item The network $\mathcal{H}_\theta$ hallucinates the missing details.
        \item The term $R(A \mathcal{H}_\theta)$ removes components that violate data consistency, acting as a soft projection into a regularised null space.
    \end{itemize}
\end{frame}

%================================================================%
\section{Variational Methods \& Optimisation Algorithms}
%================================================================%

\begin{frame}{Moving Beyond the Pseudo-Inverse}
    Because $A^\dagger$ is intractable to compute directly for large systems, classical mathematics provides a workaround: \textbf{Variational Methods}.
    
    Instead of a direct inversion, we formulate the recovery as an optimisation problem yielding
    \begin{align*}
        R_\alpha(y^\delta) = \argmin_x \left\{ \mathcal{D}_{y^\delta}(Ax) + \alpha \mathcal{J}(x) \right\} \, ,
    \end{align*}
    where
    \begin{itemize}
        \item $\mathcal{D}_{y^\delta}$ is the \textbf{Data Fidelity} term, ensuring the solution matches the measurements.
        \item $\mathcal{J}$ is the \textbf{Regularisation} term, enforcing prior knowledge about $x^\dagger$.
        \item $\alpha > 0$ is a hyper-parameter balancing the two.
    \end{itemize}
\end{frame}

\begin{frame}{Examples of Variational Terms}
    \textbf{Data Fidelity $\mathcal{D}_{y^\delta}(Ax)$:}
    \begin{itemize}
        \item $\frac{1}{2}\|Ax - y^\delta\|_2^2$: Standard Gau\ss ian noise assumption.
        \item $\|Ax - y^\delta\|_1$: Robust to Laplacian-distributed, heavy-tailed or outlier noise.
        \item Kullback--Leibler Divergence $\sum_{i} \left( y_i^\delta \log\left(\frac{y_i^\delta}{(Ax)_i}\right) + (Ax)_i - y_i^\delta \right)$: Used for Poisson noise (e.g., photon counting in PET).
    \end{itemize}
    
    \vspace{0.5em}
    \textbf{Regularisation $\mathcal{J}(x)$:}
    \begin{itemize}
        \item $\frac{1}{2}\|x\|_2^2$: Tikhonov regularisation (promotes small energy).
        \item $\|x\|_1$: Promotes sparsity (many pixels/coefficients are exactly zero).
        \item $\sup_{\| p \|_\infty \leq 1} \langle x, \operatorname{div} p \rangle \approx \|\nabla x\|_1$: Total Variation (TV) (promotes piecewise constant images, preserving sharp edges).
    \end{itemize}
\end{frame}

\begin{frame}{Convex Functions and Global Optimality}
    A critical property in variational methods is \textbf{Convexity}. 
    \begin{block}{Convex Function}
        A function $E \colon \mathcal{X} \to \mathbb{R} \cup \{+\infty\}$ is convex if for all $x, y \in \mathcal{X}$ and $\lambda \in [0, 1]$:
        \begin{align*}
            E(\lambda x + (1 - \lambda)y) \leq \lambda E(x) + (1 - \lambda) E(y) \, .
        \end{align*}
    \end{block}
    
    \begin{theoremblock}{Global Optimality}
        If $E$ is a proper, convex, and lower semi-continuous function, any local minimum is a global minimum. Furthermore, the set of global minimisers is convex.
    \end{theoremblock}
    
    \textbf{Example:} The data fidelity $\mathcal{D}_{y^\delta}(Ax) = \frac{1}{2}\|Ax - y^\delta\|_2^2$ is convex. If $\mathcal{J}(x) = \|x\|_1$ is also convex, their sum $E(x) = \mathcal{D}_{y^\delta}(Ax) + \alpha \mathcal{J}(x)$ is convex.
\end{frame}

\begin{frame}{Smooth Optimisation (Recall)}
    If $E(x)$ is continuously differentiable, a minimiser $x^*$ satisfies the \textbf{first-order optimality condition} given by
    \begin{align*}
        \nabla E(x^*) = 0 \, .
    \end{align*}
    
    As derived on Monday via the Bregman proximal framework, we can compute $x^*$ using the \textbf{Gradient Descent} iteration yielding
    \begin{align*}
        x^{k+1} = x^k - \tau \nabla E(x^k) \, .
    \end{align*}
    
    \textbf{Challenge:} Many powerful regularisation functions (like the $\ell_1$-norm or Total Variation) are convex but \textbf{non-smooth}. The gradient $\nabla E$ does not exist everywhere.
\end{frame}

\begin{frame}{Subdifferential Calculus}
    \only<1>{To handle non-smoothness, we generalise the concept of a gradient.}
    \begin{block}{Subgradient and Subdifferential}
        A vector $g$ is a \textbf{subgradient} of a convex function $E$ at $x$ if it provides a global linear lower bound satisfying
        \begin{align*}
            E(z) \geq E(x) + \langle g, z - x \rangle \quad \forall z \in \mathcal{X} \, .
        \end{align*}
        The set of all subgradients at $x$ is called the \textbf{subdifferential}, denoted $\partial E(x)$.
    \end{block}
    
    \only<2->{\begin{example}[Absolute Value]
        For $E(x) = |x|$, the subdifferential reads
        $
            \partial |x| = \begin{cases} 
            \{1\} & x > 0 \\ 
            \{-1\} & x < 0 \\ 
            [-1, 1] & x = 0 
            \end{cases}
        $
    \end{example}}
\end{frame}

\begin{frame}{Generalised Fermat's Rule}
    With the subdifferential defined, we can rigorously generalise the optimality condition.
    \begin{theoremblock}{Fermat's Rule for Convex Functions}
        A point $x^*$ is a global minimiser of a proper convex function $E$ if and only if
        \begin{align*}
            0 \in \partial E(x^*) \, .
        \end{align*}
    \end{theoremblock}\pause
    
    \textbf{Proof:} If $0 \in \partial E(x^*)$, then by definition $E(z) \geq E(x^*) + \langle 0, z - x^* \rangle = E(x^*)$ for all $z$, making $x^*$ a global minimum.\pause
    
    \vspace{0.5em}
    While we could use \textbf{Subgradient Descent} ($x^{k+1} = x^k - \tau_k g^k$ for $g^k \in \partial E(x^k)$) to find a global minimum, it requires diminishing step sizes and converges slowly ($\mathcal{O}(1/\sqrt{k})$). This motivates \textbf{Proximal Algorithms}.
\end{frame}

\begin{frame}{The Proximal Operator}
    \only<1>{To efficiently handle non-smooth convex regularisation functions, we define the \textbf{Proximal Operator}.}
    
    \begin{block}{Definition}
        For a proper, convex, lower semi-continuous function $f$ and a scalar $\tau > 0$, the proximal operator is defined via
        \begin{align*}
            \prox_{\tau f}(v) := \argmin_x \left\{ f(x) + \frac{1}{2\tau} \|x - v\|_2^2 \right\} \, .
        \end{align*}
    \end{block}
    
    \only<2->{\begin{itemize}
        \item \textbf{Uniqueness:} The quadratic penalty is strongly convex, making the minimiser strictly unique.
        \item \textbf{Optimality:} By Fermat's Rule, $x^* = \prox_{\tau f}(v)$ if and only if
        \begin{align*}
            0 \in \partial f(x^*) + \frac{1}{\tau}(x^* - v) \quad \iff \quad v \in x^* + \tau \partial f(x^*) \, .
        \end{align*}
    \end{itemize}}
\end{frame}
    
\begin{frame}{Example: Soft-Thresholding}
    Consider the $\ell_1$-norm regularisation, i.e., $f(x) = \alpha \|x\|_1$. 
    
    To find $x^* = \prox_{\tau \alpha \|\cdot\|_1}(v)$, we use the optimality condition (for each element)
    \begin{align*}
        v \in x^* + \tau \alpha \partial |x^*| \, .
    \end{align*}
    Checking the cases for the subdifferential of the absolute value yields
    \begin{itemize}
        \item If $x^* > 0$, $\partial |x^*| = \{1\} \implies v = x^* + \tau \alpha \implies x^* = v - \tau \alpha$.
        \item If $x^* < 0$, $\partial |x^*| = \{-1\} \implies v = x^* - \tau \alpha \implies x^* = v + \tau \alpha$.
        \item If $x^* = 0$, $\partial |x^*| = [-1, 1] \implies v \in [-\tau \alpha, \tau \alpha]$.
    \end{itemize}
    Combining these cases, we obtain the \textbf{Soft-Thresholding} operator yielding
    \begin{align*}
        \operatorname{soft-thresh}_{\tau\alpha}(v) = \operatorname{sign}(v) \max(|v| - \tau\alpha, 0) \, .
    \end{align*}
\end{frame}

\begin{frame}{Proximal Operators as Activation Functions}
    Many standard neural network activation functions can be interpreted as proximal operators for specific convex functions $f$, providing a natural bridge to deep learning.
    
    \vspace{0.5em}
    \begin{columns}[T]
        \column{0.33\textwidth}
        \centering
        \textbf{Soft-Thresholding} \\
        {\footnotesize $f(x) = \alpha|x| \vphantom{\frac{1}{2}\ln(1-x^2)}$}
        
        \vspace{0.5em}
        \begin{tikzpicture}[scale=0.8]
            \begin{axis}[width=5.5cm, height=5cm, axis lines=middle, xmin=-3.5, xmax=3.5, ymin=-2.5, ymax=2.5, xtick=\empty, ytick=\empty, xlabel={$x$}, ylabel={$y$}]
                \addplot[blue, thick, domain=-3:-1] {x+1};
                \addplot[blue, thick, domain=-1:1] {0};
                \addplot[blue, thick, domain=1:3] {x-1};
                \addplot[dashed, gray] coordinates {(-2.5,-2.5) (2.5,2.5)};
            \end{axis}
        \end{tikzpicture}
        
        \column{0.33\textwidth}
        \centering
        \textbf{ReLU} \\
        {\footnotesize $f(x) = \chi_{[0, \infty)}(x) \vphantom{\frac{1}{2}\ln(1-x^2)}$}
        
        \vspace{0.5em}
        \begin{tikzpicture}[scale=0.8]
            \begin{axis}[width=5.5cm, height=5cm, axis lines=middle, xmin=-3.5, xmax=3.5, ymin=-2.5, ymax=2.5, xtick=\empty, ytick=\empty, xlabel={$x$}, ylabel={$y$}]
                \addplot[red, thick, domain=-3:0] {0};
                \addplot[red, thick, domain=0:3] {x};
                \addplot[dashed, gray] coordinates {(-2.5,-2.5) (2.5,2.5)};
            \end{axis}
        \end{tikzpicture}
        
        \column{0.33\textwidth}
        \centering
        \textbf{Tanh} \\
        {\footnotesize $f(x) = x \operatorname{arctanh}(x) + \frac{1}{2}\ln(1-x^2) - \frac{1}{2}x^2$}
        
        \vspace{0.5em}
        \begin{tikzpicture}[scale=0.8]
            \begin{axis}[width=5.5cm, height=5cm, axis lines=middle, xmin=-3.5, xmax=3.5, ymin=-2.5, ymax=2.5, xtick=\empty, ytick=\empty, xlabel={$x$}, ylabel={$y$}]
                \addplot[darkgreen, thick, domain=-3:3, samples=100] {tanh(x)};
                \addplot[dashed, gray] coordinates {(-2.5,-2.5) (2.5,2.5)};
            \end{axis}
        \end{tikzpicture}
    \end{columns}
\end{frame}

\begin{frame}{Proximal Gradient Descent (ISTA)}
    We can now solve problems of the form $\min_x \mathcal{D}_{y^\delta}(Ax) + \alpha \mathcal{J}(x)$, where $\mathcal{D}_{y^\delta}$ is smooth but $\mathcal{J}$ is non-smooth.
    
    The \textbf{Iterative Shrinkage-Thresholding Algorithm (ISTA)} alternates between a gradient descent step on the data fidelity, and a proximal step on the regulariser.
    
    \begin{theoremblock}{ISTA Algorithm}
        For $k = 0, 1, 2, \dots$ until convergence we iterate the updates
        \begin{align*}
            z^k &= x^k - \eta A^\top \nabla \mathcal{D}_{y^\delta}(Ax^k) = x^k - \eta A^\top(Ax^k - y^\delta) \\
            x^{k+1} &= \prox_{\eta \alpha \mathcal{J}}(z^k) := \argmin_x \left\{ \frac{1}{2\eta} \|x - z^k\|^2 + \alpha \mathcal{J}(x) \right\} \, .
        \end{align*}
    \end{theoremblock}
\end{frame}

\begin{frame}{Bregman Proximal View of ISTA}
    \only<1-2>{Similar to Gradient Descent, ISTA can be elegantly derived from the \textbf{Bregman proximal point method}.}
    
    Consider the iterative update given by
    \begin{align*}
        x^{k+1} = \argmin_x \left\{ \mathcal{D}_{y^\delta}(Ax) + \alpha \mathcal{J}(x) + D_H(x, x^k) \right\} \, ,
    \end{align*}
    where $D_H$ is the Bregman distance.\pause
    
    If we choose the distance generating function $H(x) = \frac{1}{2\eta}\|x\|_2^2 - \mathcal{D}_{y^\delta}(Ax)$, the data fidelity is locally linearised. Expanding the objective yields the quadratic model
    \begin{align*}
        x^{k+1} &= \argmin_x \left\{ \mathcal{D}_{y^\delta}(Ax^k) + \langle A^\top \nabla \mathcal{D}_{y^\delta}(Ax^k), x - x^k \rangle + \frac{1}{2\eta}\|x - x^k\|_2^2 + \alpha \mathcal{J}(x) \right\} \\
        &= \argmin_x \left\{ \frac{1}{2\eta} \left\| x - \left( x^k - \eta A^\top \nabla \mathcal{D}_{y^\delta}(Ax^k) \right) \right\|_2^2 + \alpha \mathcal{J}(x) \right\} \, .
    \end{align*}
    \only<3->{This exactly recovers the ISTA update! For $H$ to be convex, we require the step size to satisfy $\eta \leq 1/\beta$, where $\beta$ is the Lipschitz constant of $A^\top \nabla \mathcal{D}_{y^\delta}(Ax)$.}
\end{frame}

\begin{frame}{Accelerating ISTA: FISTA}
    While ISTA is mathematically sound, it converges at a slow sublinear rate of $\mathcal{O}(1/k)$.\pause
    
        The \textbf{Fast Iterative Shrinkage-Thresholding Algorithm (FISTA)}~\citep{Beck2009FISTA} accelerates this to an optimal $\mathcal{O}(1/k^2)$ rate by introducing \textbf{Nesterov momentum}~\citep{Nesterov1983Method}.\pause
    
    Instead of taking the gradient step from the current point $x^k$, FISTA evaluates the gradient at an extrapolated point $y^k$, yielding
    \begin{align*}
        y^k &= x^k + \left( \frac{t_{k-1} - 1}{t_k} \right) (x^k - x^{k-1}) \, , \\
        x^{k+1} &= \prox_{\eta \alpha \mathcal{J}}\left( y^k - \eta A^\top(Ay^k - y^\delta) \right) \, ,
    \end{align*}
        where $t_k$ is a deterministic momentum sequence. \pause This significantly speeds up convergence with negligible extra computational cost, and convergence of the iterates themselves can be guaranteed under specific choices of this sequence~\citep{chambolle2015convergence}.
\end{frame}

\begin{frame}{Limitations of ISTA}
    ISTA is simple and elegant, but it requires the data fidelity gradient to be Lipschitz continuous.
    
    \begin{alertblock}{The Problem with $A^\top A$}
        The gradient step requires evaluating $\eta A^\top A x^k$. If the operator $A$ is highly ill-conditioned, $A^\top A$ has extreme eigenvalues. 
        To ensure convergence, the step size $\eta$ must be strictly less than $1/L$ (where $L = \|A^\top A\|$). If $L$ is huge, $\eta$ becomes tiny, and ISTA takes thousands of iterations to converge.
    \end{alertblock}\pause
    
    Furthermore, what if the regulariser involves a linear operator inside the norm, like Total Variation $\|\nabla x\|_1$? We do not have a closed-form proximal operator for this.
\end{frame}

\begin{frame}{Primal--Dual Hybrid Gradient (PDHG)}
    \only<1-3>{To avoid the explicit evaluation of $A^\top A$ and handle more complex regularisation functions, we can use \textbf{Primal--Dual Algorithms} (e.g., Chambolle--Pock~\citep{Pock2009PDHG, Chambolle2011PDHG}).}\pause
    
    These methods introduce a \textbf{dual variable} $h$ that lives in the measurement space, allowing us to decouple the operators.\pause
    
    To solve $\min_x \mathcal{J}(x) + \mathcal{D}_{y^\delta}(Ax)$, we alternate the updates yielding
    \begin{align*}
        x^{k+1} &= \prox_{\tau \mathcal{J}}\left( x^k - \tau A^\top h^k \right) \quad &\text{(Primal Update in } \mathcal{X} \text{)} \\
        \bar{x}^{k+1} &= x^{k+1} + \omega(x^{k+1} - x^k) \quad &\text{(Extrapolation)} \\
        h^{k+1} &= \prox_{\sigma \mathcal{D}_{y^\delta}^*}\left( h^k + \sigma A \bar{x}^{k+1} \right) \quad &\text{(Dual Update in measurement space)} \, ,
    \end{align*}
    where $\mathcal{D}_{y^\delta}^*$ is the convex conjugate of the data fidelity, and $\omega$ is often set to $1$, simplifying the extrapolation step to $2x^{k+1} - x^k$.
    
    \only<4->{Notice we only apply $A$ and $A^\top$ individually; they are never multiplied together!}
\end{frame}

\begin{frame}{Connection to the Coding Challenge: Condat-V\~u}
    \only<1-3>{In your \textbf{Coding Challenge}, the baseling algorithm is the \textbf{Condat-V\~u} algorithm~\citep{condat2013primal, vu2013splitting}.}\pause
    
    Condat-V\~u is an extension of Primal--Dual splitting designed to efficiently handle an objective with \emph{three} parts yielding
    \begin{align*}
        \min_x \underbrace{\mathcal{D}_{y^\delta}(Ax)}_{\text{Smooth}} + \underbrace{G(x)}_{\text{Proximable}} + \underbrace{\mathcal{H}(Dx)}_{\text{Non-smooth composition}} \, .
    \end{align*}\pause
    For example, setting $G(x) = 0$ and $\mathcal{H}(Dx) = \alpha \|Dx\|_1$ recovers Total Variation regularisation where $\mathcal{J}(x) = \mathcal{H}(Dx)$ and we split the finite difference operator $D$.\pause
    
    \only<4->{The updates resemble PDHG but include the additional derivative of the smooth term yielding
    \begin{align*}
        x^{k+1} &= \prox_{\tau G}\left( x^k - \tau A^\top \nabla \mathcal{D}_{y^\delta}(Ax^k) - \tau D^\top h^k \right) \, , \quad &\text{(Primal Update)} \\
        \bar{x}^{k+1} &= 2x^{k+1} - x^k \, , \quad &\text{(Extrapolation)} \\
        h^{k+1} &= \prox_{\sigma \mathcal{H}^*}\left( h^k + \sigma D \bar{x}^{k+1} \right) \, . \quad &\text{(Dual Update)} 
    \end{align*}
    \textbf{Caution:} By treating the data fidelity as a smooth term, we require the Lipschitz constant of its derivative, meaning we must compute $\|A^\top A\|$ again!}
\end{frame}

%================================================================%
\section{Bilevel Optimisation \& Sensitivity Analysis}
%================================================================%

\begin{frame}{The hyperparameter problem}
    In standard machine learning, we minimise an empirical risk on a training set to find model parameters. However, this process is governed by \textbf{hyperparameters} $\Theta$ (e.g., regularisation weights), which control the model's capacity and generalisation.
    
    \vspace{0.5em}
    Typically, we select $\Theta$ by evaluating the trained model on a separate \textbf{validation set}. As the dimensionality of $\Theta$ grows, manual tuning or grid search becomes intractable.
    
    \vspace{0.5em}
    Similarly, classical variational methods for inverse problems are mathematically robust but rely on manual tuning of regularisation parameters (like the weight $\alpha$). If $\alpha$ is too small, the image remains noisy; if too large, it loses critical features.
\end{frame}

\begin{frame}{Bilevel optimisation framework}
    We can formalise hyperparameter tuning elegantly as a nested \textbf{bilevel optimisation} problem. For a validation set $\mathcal{S}_{\text{val}}$ and training set $\mathcal{S}_{\text{train}}$, we observe
    \begin{align*}
        \min_{\Theta} \; \frac{1}{|\mathcal{S}_{\text{val}}|} \sum_{(x,y) \in \mathcal{S}_{\text{val}}} \ell_{\text{val}}\!\left( f_{w^\star(\Theta)}(x), y \right)
    \end{align*}
    subject to the lower-level training problem yielding
    \begin{align*}
        w^\star(\Theta) \in \argmin_w \; \frac{1}{|\mathcal{S}_{\text{train}}|} \sum_{(x,y) \in \mathcal{S}_{\text{train}}} \ell_{\text{train}}\!\left( f_w(x), y; \Theta \right) \, .
    \end{align*}
    \begin{itemize}
        \item \textbf{Lower level:} Fit the model parameters $w$ to the training data for fixed hyperparameters $\Theta$.
        \item \textbf{Upper level:} Update $\Theta$ to minimise the generalisation error on the validation data.
    \end{itemize}
\end{frame}

\begin{frame}{Example: tuning ridge regression}
    Consider learning the optimal penalty $\Theta$ for a linear model $f_w(x) = x^\top w$. The bilevel formulation yields
    \begin{itemize}
        \item \textbf{Lower level (Training):} Compute regularised weights
        \begin{align*}
            w^\star(\Theta) = \argmin_w \; \frac{1}{2|\mathcal{S}_{\text{train}}|} \sum_{(x,y) \in \mathcal{S}_{\text{train}}} \left( x^\top w - y \right)^2 + \frac{\Theta}{2}\|w\|^2 \, .
        \end{align*}
        \item \textbf{Upper level (Validation):} Minimise generalisation error
        \begin{align*}
            \min_{\Theta > 0} \; \frac{1}{2|\mathcal{S}_{\text{val}}|} \sum_{(x,y) \in \mathcal{S}_{\text{val}}} \left( x^\top w^\star(\Theta) - y \right)^2 \, .
        \end{align*}
    \end{itemize}
    \vspace{0.5em}
    Optimising $\Theta$ on a hold-out validation set prevents the trivial, overfitted solution $\Theta = 0$ that would occur if we minimised the upper-level loss strictly on the training set.
\end{frame}

\begin{frame}{Bilevel learning for inverse problems}
    Applying this to inverse problems, suppose we have a dataset of ground-truth images $x_i^\dagger$ and noisy measurements $y_i^\delta$. We can \emph{learn} the optimal parameters $\Theta$ of our regulariser yielding
    \begin{align*}
        \min_{\Theta} \frac{1}{N} \sum_{i=1}^N \mathcal{L}\left( x_i^\star(\Theta), \, x_i^\dagger \right) \quad \text{subject to} \quad x_i^\star(\Theta) \in \argmin_x E(x, \Theta; y_i^\delta) \, ,
    \end{align*}
    where $E(x, \Theta; y_i^\delta) := \mathcal{D}_{y_i^\delta}(Ax) + \mathcal{R}_\Theta(x)$.
    
    \begin{itemize}
        \item \textbf{Lower level:} Solve the variational inverse problem to find the optimal reconstructed image $x^\star$ for fixed parameters $\Theta$.
        \item \textbf{Upper level:} Evaluate the reconstruction quality against the ground truth $x^\dagger$, and update $\Theta$.
    \end{itemize}
\end{frame}

\begin{frame}{Example: tuning Tikhonov regularisation}
    For a linear inverse problem $y^\delta = Ax + \eta$, tuning the Tikhonov parameter $\Theta$ uses a dataset of ground-truth signals $x_i^\dagger$. This yields
    \begin{itemize}
        \item \textbf{Lower level (Reconstruction):} Solve the regularised problem
        \begin{align*}
            x_i^\star(\Theta) = \argmin_x \; \frac{1}{2} \|Ax - y_i^\delta\|^2 + \frac{\Theta}{2}\|x\|^2 \, .
        \end{align*}
        \item \textbf{Upper level (Supervised Training):} Minimise reconstruction error
        \begin{align*}
            \min_{\Theta > 0} \; \frac{1}{2N} \sum_{i=1}^N \left\| x_i^\star(\Theta) - x_i^\dagger \right\|^2 \, .
        \end{align*}
    \end{itemize}
    \vspace{0.5em}
    The lower-level problem is mathematically identical to ridge regression. The distinction lies entirely in the upper level, which measures distance to the exact ground truth rather than predictive risk.
\end{frame}

\begin{frame}{The gradient challenge}
    To perform gradient descent on the upper level, we need to compute the gradient of the loss with respect to the hyperparameters $\Theta$ given by
    \begin{align*}
        \nabla_\Theta \mathcal{L} = \nabla_{x^\star} \mathcal{L} \frac{\partial x^\star}{\partial \Theta} \, .
    \end{align*}
    
    The term $\nabla_{x^\star} \mathcal{L}$ is easy (e.g., the derivative of the mean squared error). 
    
    \textbf{The challenge:} How do you compute $\frac{\partial x^\star}{\partial \Theta}$? This requires differentiating through the $\argmin$ operation in the lower level!
    
    \vspace{0.5em}
    We cannot simply use automatic differentiation (e.g., PyTorch's \texttt{autograd}) to backpropagate through thousands of ISTA iterations, as it would require storing the entire computational graph in memory.
\end{frame}

\begin{frame}{The implicit function theorem (IFT)}
    Instead of tracking the iterations, we use the \textbf{implicit function theorem}. 
    
    Since $x^\star$ is the minimum of the lower-level energy $E$, it must satisfy the first-order optimality condition yielding
    \begin{align*}
        \nabla_x E(x^\star, \Theta) = 0 \, .
    \end{align*}
    This equation implicitly defines $x^\star$ as a continuous function of $\Theta$.
\end{frame}

\begin{frame}{Deriving the sensitivity via IFT}
    We take the total derivative of the optimality condition $\nabla_x E(x^\star, \Theta) = 0$ with respect to $\Theta$. Applying the chain rule yields
    \begin{align*}
        \nabla_{xx}^2 E(x^\star, \Theta) \frac{\partial x^\star}{\partial \Theta} + \nabla_{\Theta x}^2 E(x^\star, \Theta) = 0 \, ,
    \end{align*}
    where $\nabla_{xx}^2 E$ is the Hessian matrix with respect to $x$, and $\nabla_{\Theta x}^2 E$ is the mixed partial derivative.
    
    \vspace{0.5em}
    Rearranging gives the exact sensitivity yielding
    \begin{align*}
        \frac{\partial x^\star}{\partial \Theta} = - \left( \nabla_{xx}^2 E(x^\star, \Theta) \right)^{-1} \nabla_{\Theta x}^2 E(x^\star, \Theta) \, .
    \end{align*}
    This is a massive breakthrough! We can analytically compute gradients for hyperparameter learning by solving a single linear system at the end, regardless of how many iterations it took to find $x^\star$.
\end{frame}

\begin{frame}{The adjoint system}
    In practice, we never invert the huge Hessian matrix. We define the \textbf{adjoint state} $p$ as the solution to the linear system satisfying
    \begin{align*}
        \left( \nabla_{xx}^2 E(x^\star, \Theta) \right) p = \nabla_{x^\star} \mathcal{L} \, .
    \end{align*}
    
    Since the Hessian is symmetric, we can substitute this into our gradient expression to obtain
    \begin{align*}
        \nabla_\Theta \mathcal{L} = - \nabla_{\Theta x}^2 E(x^\star, \Theta) \, p \, .
    \end{align*}
    
    We solve for $p$ using an iterative solver (like Conjugate Gradient), and then compute the final hypergradient efficiently.
\end{frame}

\begin{frame}{Recall from Monday: implicit layer backpropagation}
    Does this adjoint system look familiar? 
    
    \vspace{0.5em}
    Yesterday, we discussed \textbf{Case I: Bilevel or optimisation-defined layers} as a generalisation of backpropagation yielding
    \begin{align*}
        z^\star(w,x) \in \argmin_{z} E(z,w;x) \, .
    \end{align*}
    We identified the constraint as $\mathcal{G}(z,w,x) = \nabla_z E(z,w;x) = 0$.
    
    \vspace{0.5em}
    The adjoint equation we derived today for the hypergradient using the IFT is \emph{exactly} the continuous Lagrangian stationarity condition from yesterday! The implicit function theorem is simply the mathematical engine driving implicit differentiation.
\end{frame}

\begin{frame}{IFT for ridge regression: computing derivatives}
    Let us apply this to the ridge regression example. The lower-level energy is
    \begin{align*}
        E(w, \Theta) = \frac{1}{2|\mathcal{S}_{\text{train}}|} \sum_{(x,y) \in \mathcal{S}_{\text{train}}} \left( x^\top w - y \right)^2 + \frac{\Theta}{2}\|w\|^2 \, .
    \end{align*}
    The first-order optimality condition for $w^\star$ yields
    \begin{align*}
        \nabla_w E(w^\star, \Theta) = \frac{1}{|\mathcal{S}_{\text{train}}|} \sum_{(x,y) \in \mathcal{S}_{\text{train}}} x \left( x^\top w^\star - y \right) + \Theta w^\star = 0 \, .
    \end{align*}\pause
    We compute the required second-order derivatives evaluated at $w^\star$:
    \begin{itemize}
        \item \textbf{Hessian (w.r.t.\ $w$):} $\nabla_{ww}^2 E(w^\star, \Theta) = \frac{1}{|\mathcal{S}_{\text{train}}|} \sum_{(x,y) \in \mathcal{S}_{\text{train}}} x x^\top + \Theta I$.
        \item \textbf{Mixed derivative (w.r.t.\ $\Theta$ and $w$):} $\nabla_{\Theta w}^2 E(w^\star, \Theta) = w^\star$.
    \end{itemize}
\end{frame}

\begin{frame}{IFT for ridge regression: sensitivity and gradient}
    The upper-level objective is $\mathcal{L}(\Theta) := \frac{1}{2|\mathcal{S}_{\text{val}}|} \sum_{(x,y) \in \mathcal{S}_{\text{val}}} \left( x^\top w^\star(\Theta) - y \right)^2$.\pause
    
    Using the chain rule, the hypergradient is $\nabla_\Theta \mathcal{L}(\Theta) = \left\langle \nabla_w \mathcal{L}(w^\star), \frac{\partial w^\star}{\partial \Theta} \right\rangle$. Substituting our derivatives yields
    \begin{align*}
        \nabla_\Theta \mathcal{L}(\Theta) = \left\langle \frac{1}{|\mathcal{S}_{\text{val}}|} \sum_{(x,y) \in \mathcal{S}_{\text{val}}} x \left( x^\top w^\star - y \right), \, - \left( \frac{1}{|\mathcal{S}_{\text{train}}|} \sum_{(x,y) \in \mathcal{S}_{\text{train}}} x x^\top + \Theta I \right)^{-1} w^\star \right\rangle \, .
    \end{align*}\pause
    
    \vspace{-0.15cm}
    Since the Hessian is symmetric and positive definite for $\Theta > 0$, we can move its inverse to the left side of the inner product to obtain the closed-form gradient yielding
    \begin{align*}
        \nabla_\Theta \mathcal{L}(\Theta) = - \left\langle \underbrace{\left( \frac{1}{|\mathcal{S}_{\text{train}}|} \sum_{(x,y) \in \mathcal{S}_{\text{train}}} x x^\top + \Theta I \right)^{-1} \frac{1}{|\mathcal{S}_{\text{val}}|} \sum_{(x,y) \in \mathcal{S}_{\text{val}}} x \left( x^\top w^\star - y \right)}_{\text{adjoint state } p}, \, w^\star \right\rangle \, .
    \end{align*}
\end{frame}

\begin{frame}{Optimising the hyperparameter}
    With the hypergradient $\nabla_\Theta \mathcal{L}(\Theta)$ computed, we can find the optimal penalty $\Theta$ using gradient descent.\pause
    
    \vspace{0.5em}
    However, we must enforce the constraint $\Theta \geq 0$. We achieve this via \textbf{projected gradient descent} yielding
    \begin{align*}
        \Theta^{k+1} = \max\left( 0, \, \Theta^k - \eta \, \nabla_\Theta \mathcal{L}(\Theta^k) \right) \, .
    \end{align*}\pause
    
    \begin{block}{Connection to proximal algorithms}
        Recall that the projection onto the non-negative orthant $\mathbb{R}_{\geq 0}$ is exactly the proximal operator for the characteristic function $\chi_{\mathbb{R}_{\geq 0}}$.
        
        \vspace{0.5em}
        Therefore, this update is a \textbf{proximal gradient descent} step. This connects our bilevel hyperparameter learning back to the proximal methods (like ISTA) we introduced earlier.
    \end{block}
\end{frame}

\begin{frame}{The issue of non-smoothness}
    \begin{alertblock}{The Hessian problem}
        The IFT formulation strictly requires the energy function $E$ to be twice differentiable, because we need the Hessian $\nabla_{xx}^2 E$.
    \end{alertblock}
    
    If our regularisation is the $\ell_1$-norm or Total Variation, it contains an absolute value function $|x|$, which is non-smooth.
    \begin{itemize}
        \item The first derivative is the sign function (undefined at $0$).
        \item The second derivative (Hessian) is $0$ almost everywhere, and undefined at $0$.
    \end{itemize}
    The IFT matrix inversion will fail if we naively apply it here.
\end{frame}

\begin{frame}{Huber smoothing}
    To use IFT with sparsity-promoting regularisation functions, we can use \textbf{Huber Smoothing}. We replace the sharp absolute value $|x|$ with a smooth, piecewise approximation given by
    \begin{align*}
        h_\epsilon(x) = 
        \begin{cases} 
            \frac{x^2}{2\epsilon} & \text{if } |x| \le \epsilon \\ 
            |x| - \frac{\epsilon}{2} & \text{if } |x| > \epsilon
        \end{cases}
        \, .
    \end{align*}
    
    \begin{columns}
    \column{0.6\textwidth}
    \begin{itemize}
        \item Near zero (inside $[-\epsilon, \epsilon]$), it is a quadratic function.
        \item Far from zero, it behaves like the $\ell_1$-norm.
        \item It is everywhere continuously twice-differentiable, allowing the Hessian to be computed safely.
    \end{itemize}
    \column{0.4\textwidth}
    \begin{tikzpicture}[scale=0.6]
        \begin{axis}[
            axis lines=middle,
            xmin=-3, xmax=3, ymin=0, ymax=3,
            xlabel=$x$, ylabel=$y$,
            legend style={at={(0.5,-0.2)},anchor=north}
        ]
        \addplot[red, thick, domain=-3:3, samples=100] {abs(x)};
        \addplot[blue, thick, dashed, domain=-3:3, samples=100] {abs(x)<=1 ? x^2/(2*1) : abs(x) - 1/2};
        \legend{$|x|$, $h_\epsilon(x)$}
        \end{axis}
    \end{tikzpicture}
    \end{columns}
\end{frame}

%================================================================%
\section{The Unrolling Principle \& Architectures}
%================================================================%

\begin{frame}{The Limits of Implicit Differentiation}
    While bilevel optimisation with the IFT is mathematically elegant, it has practical drawbacks:
    \begin{itemize}
        \item \textbf{Speed:} The lower-level solver must be run to complete convergence (often thousands of iterations) for every single training example.
        \item \textbf{Adjoint Cost:} Solving the adjoint linear system adds significant overhead.
        \item \textbf{Restricted Expressivity:} We are only learning a few scalar hyperparameters (like $\alpha$), not full convolutional filters.
    \end{itemize}
\end{frame}

\begin{frame}{Truncating the Solver}
    What if, instead of waiting for convergence $k \to \infty$, we simply truncate the iterative algorithm (like ISTA) after a fixed, small number of steps $K$ (e.g., $K=10$)?
    
    If $K$ is small, we do not need the IFT. We can simply store the computational graph in memory and use standard PyTorch/TensorFlow backpropagation!
\end{frame}

\begin{frame}{The Unrolling Principle}
    This leads to one of the most powerful concepts in modern computational imaging: \textbf{Algorithm Unrolling} (or Unfolding).
    
    \begin{theoremblock}{The Unrolling Principle}
        Truncate an iterative algorithm to a finite number of steps $K$. Treat each iteration as a layer in a Deep Neural Network. Relax the parameters of the algorithm to be iteration-dependent, and learn them end-to-end from data using backpropagation.
    \end{theoremblock}
    
    Rather than hoping an algorithm converges in 1000 steps, we explicitly train it to produce the best possible image in exactly 10 steps.
\end{frame}

\begin{frame}{From Iterations to Layers}
    \begin{center}
    \begin{tikzpicture}[
        box/.style={draw, thick, fill=stone, minimum width=2cm, minimum height=1cm, align=center},
        net/.style={draw, thick, fill=uclblue!20, minimum width=2cm, minimum height=1cm, align=center},
        arr/.style={->, thick, >=latex}
    ]
        \node (x0) at (0,0) {$x^0$};
        \node[box] (grad1) at (2.5,0) {Iteration 1};
        \node[net] (prox1) at (5.5,0) {Iteration 2};
        \node[box] (grad2) at (8.5,0) {Iteration 3};
        \node (dots) at (11.5,0) {$\dots$};
        \node[net] (out) at (13.5,0) {Iteration $K$};
        
        \draw[arr] (x0) -- (grad1);
        \draw[arr] (grad1) -- node[above] {$x^1$} (prox1);
        \draw[arr] (prox1) -- node[above] {$x^2$} (grad2);
        \draw[arr] (grad2) -- node[above] {$x^3$} (dots);
        \draw[arr] (dots) -- (out);
        \draw[arr] (out) -- node[above] {$x^K$} (15,0);
    \end{tikzpicture}
    \end{center}
    Every block represents one step of the classical mathematical algorithm, but parameterised by neural network weights $\theta_k$.
\end{frame}

\begin{frame}{Unrolling ISTA}
    Applying unrolling to ISTA gives us \textbf{LISTA}~\citep{gregor2010learning}, originally proposed for sparse coding.
    
    Recall a standard ISTA step yielding
    \begin{align*}
        x^{k+1} = \operatorname{soft-thresh}_{\alpha\eta} \left( x^k - \eta A^\top(Ax^k - y^\delta) \right) \, .
    \end{align*}
    
    We can algebraically expand the argument of the soft-thresholding function to observe that
    \begin{align*}
        x^k - \eta A^\top A x^k + \eta A^\top y^\delta = (I - \eta A^\top A) x^k + (\eta A^\top) y^\delta \, .
    \end{align*}
\end{frame}

\begin{frame}{LISTA (Learned ISTA)}
    By defining weight matrices $W_1 = I - \eta A^\top A$ and $W_2 = \eta A^\top$, the ISTA step looks exactly like a neural network layer satisfying
    \begin{align*}
        x^{k+1} = \operatorname{soft-thresh}_{\theta_k} \left( W_1 x^k + W_2 y^\delta \right) \, ,
    \end{align*}
    where the soft-thresholding operator acts as the non-linear activation function (similar to ReLU)!
    
    \vspace{0.5em}
    In LISTA, we stop enforcing that $W_1$ and $W_2$ are strictly derived from the physics operator $A$. Instead, we make them \textbf{learnable weight matrices} $W_1^k$, $W_2^k$, unique to each layer $k$.
\end{frame}

\begin{frame}{Decoupling Parameters in LISTA}
    By decoupling the matrices from the forward operator, LISTA gains immense flexibility:
    \begin{itemize}
        \item The network can learn matrices that implicitly condition the problem, accelerating convergence.
        \item The threshold parameter $\theta_k$ is learned independently for each layer, allowing the network to adaptively relax regularisation as the image gets closer to the ground truth.
    \end{itemize}
    
    \textbf{Result:} LISTA achieves in $K=10$ iterations what standard ISTA achieves in $K=1000$ iterations!
\end{frame}

\begin{frame}{Accelerating convergence with LISTA}
    By learning a data-driven preconditioning, LISTA achieves a near-linear convergence rate within the trained horizon, drastically outperforming classical analytical methods.
    
    \vspace{0.5em}
    \centering
    \begin{tikzpicture}[scale=0.8]
        \begin{axis}[
            xlabel={Iterations $k$},
            ylabel={Error (log scale)},
            ymode=log,
            domain=1:20,
            legend pos=north east,
            grid=major,
            width=10cm, height=5.5cm
        ]
        \addplot[blue, thick] {1/x}; \addlegendentry{ISTA $\mathcal{O}(1/k)$}
        \addplot[red, thick] {1/(x^2)}; \addlegendentry{FISTA $\mathcal{O}(1/k^2)$}
        \addplot[darkgreen, thick, mark=*] coordinates {
            (1, 1.0) (2, 0.1) (3, 0.01) (4, 0.001) (5, 0.0001)
        }; \addlegendentry{LISTA (Trained)}
        \end{axis}
    \end{tikzpicture}
\end{frame}

\begin{frame}{Unrolling PDHG}
    LISTA is great for sparse coding, but for complex imaging inverse problems (like CT or MRI), the matrices $W_1$ and $W_2$ would be impractically large to learn. 
    
    Instead, we unroll the Primal--Dual Hybrid Gradient (PDHG) algorithm.
    
    Recall that the PDHG updates read
    \begin{align*}
        x^{k+1} &= \prox_{\tau \mathcal{J}}\left( x^k - \tau A^\top h^k \right) \, , \\
        h^{k+1} &= \prox_{\sigma \mathcal{D}_{y^\delta}^*}\left( h^k + \sigma A \bar{x}^{k+1} \right) \, .
    \end{align*}
        
        \begin{alertblock}{The step-size bottleneck}
            Analytical PDHG guarantees convergence only if the step sizes satisfy $\tau \sigma \|A\|^2 < 1$. This strict Lipschitz bound forces the algorithm to take small, conservative steps. By replacing these updates with neural networks, we can learn more aggressive, data-driven step sizes.
        \end{alertblock}
\end{frame}

\begin{frame}{Learned Primal--Dual (LPD)}
    To create the \textbf{Learned Primal--Dual (LPD)} network~\citep{adler2018learned}, we replace the analytical proximal operators with Convolutional Neural Networks (CNNs).
    
    \begin{theoremblock}{LPD Layer Equations}
        \begin{align*}
            \mathbf{h}^{k+1} &= \Gamma_{\theta_k}\!\left(\mathbf{h}^k, \, A (\mathbf{x}^k)_1, \, y^\delta\right) \quad &\text{(Dual space CNN)} \, , \\
            \mathbf{x}^{k+1} &= \Lambda_{\phi_k}\!\left(\mathbf{x}^k, \, A^\top (\mathbf{h}^{k+1})_1\right) \quad &\text{(Primal space CNN)} \, .
        \end{align*}
    \end{theoremblock}
    
    \textbf{Memory (Hidden States):} The variables $\mathbf{x}^k$ and $\mathbf{h}^k$ are expanded to have multiple channels. Only the first channel $(\cdot)_1$ represents the physical image or data; the rest act as memory passing hidden states across iterations, similar to a Recurrent Neural Network.
\end{frame}

\begin{frame}{Learned Primal--Dual (LPD) architecture diagram}
    \begin{center}
    \resizebox{0.9\textwidth}{!}{
    \begin{tikzpicture}[
        node distance=2.5cm and 2cm,
        >=stealth,
        thick,
        state/.style={rectangle, draw=blue!60!black, fill=blue!5, minimum height=1cm, minimum width=1.2cm, rounded corners=2pt, font=\bfseries},
        update/.style={rectangle, draw=red!60!black, fill=red!5, minimum height=1.2cm, minimum width=1.8cm, rounded corners=2pt, align=center, font=\small},
        operator/.style={circle, draw=green!60!black, fill=green!5, minimum size=1cm, font=\itshape},
        data/.style={rectangle, draw=gray!80, fill=gray!10, minimum size=0.8cm, rounded corners},
        skip line/.style={->, draw=black!70, dashed}
    ]
        
        \node[state] (hk) {$\mathbf{h}^k$};
        \node[state, below=3cm of hk] (uk) {$\mathbf{x}^k$};
        
        \node[update, right=2.5cm of hk] (Gamma) {$\Gamma_{\theta_k}$};
        \node[operator] (K) at ($(hk)!0.5!(Gamma) + (0, -1.5)$) {$A$};
        \node[data, above=1cm of Gamma] (f) {$y^\delta$};
        
        \node[state, right=2.5cm of Gamma] (hk1) {$\mathbf{h}^{k+1}$};
        
        \node[update] (Lambda) at (hk1 |- uk) {$\Lambda_{\phi_k}$};
        \node[operator] (Kadj) at ($(hk1)!0.5!(Lambda) + (0, 0)$) {$A^\top$};
        
        \node[state, right=2.5cm of Lambda] (uk1) {$\mathbf{x}^{k+1}$};
        
        \draw[->] (uk) -| node[pos=0.2, above, font=\tiny] {$(1)$} (K);
        \draw[->] (K) -- node[midway, right, font=\tiny] {$\text{eval}_x$} (Gamma.south west);
        \draw[->] (hk) -- (Gamma);
        \draw[->] (f) -- (Gamma);
        \draw[->] (Gamma) -- (hk1);
        
        \draw[->] (hk1) -- node[midway, right, font=\tiny] {$(1)$} (Kadj);
        \draw[->] (Kadj) -- node[midway, right, font=\tiny] {$\text{eval}_h$} (Lambda.north west);
        
        \draw[skip line] (uk.south) -- ++(0,-0.6) -| ($(Lambda.south west)+(-0.3,0)$) -- (Lambda.south west);
        \draw[->] (Lambda) -- (uk1);
        
        \node[anchor=east, blue!50!black] at ($(hk.west)+(-0.5,0)$) {Dual Space};
        \node[anchor=east, blue!50!black] at ($(uk.west)+(-0.5,0)$) {Primal Space};
    \end{tikzpicture}
    }
    \end{center}
    \vspace{0.5em}
    The CNNs often use a \textbf{UNet architecture} to capture multi-scale features, explicitly replacing the classical proximal operator. The skip connections allow state vectors to bypass iterations.
\end{frame}

\begin{frame}{Physics-Informed Learning}
    \textbf{Why is LPD so successful?}
    
    Unlike end-to-end networks that try to guess the inverse, LPD explicitly uses the forward operator $A$ and the adjoint operator $A^\top$ in every single layer.
    
    \begin{itemize}
        \item The Dual CNN $\Gamma$ learns to correct measurement noise and interpolate missing data in the raw sensor space.
        \item The operator $A^\top$ mathematically back-projects this corrected data into the image space.
        \item The Primal CNN $\Lambda$ learns to remove the specific back-projection artefacts and enhance image textures.
    \end{itemize}
    The physics ($A, A^\top$) and the data-driven priors ($\Gamma, \Lambda$) work in perfect synergy.
\end{frame}

%================================================================%
\section{Looking Ahead}
%================================================================%

\begin{frame}{Summary}
    Today we covered:
    \begin{enumerate}
        \item \textbf{Null-Space Networks:} Structurally forcing neural networks to obey measured data physics.
        \item \textbf{Variational Methods:} Formulating inverse problems mathematically with fidelity and regularisation, and solving them with ISTA and PDHG.
        \item \textbf{Bilevel Optimisation:} Learning hyper-parameters using the Implicit Function Theorem and Huber smoothing.
        \item \textbf{Algorithm Unrolling:} Truncating iterative solvers (LISTA, LPD) and replacing proximal steps with CNNs for fast, physics-informed inference.
    \end{enumerate}
\end{frame}

\begin{frame}{Preview: Wednesday and Thursday}
    Unrolling is powerful, but it ties the trained network directly to the specific forward model $A$. If the physics change slightly (e.g., a different CT scanner), you have to retrain the LPD network from scratch!
    
    \vspace{0.5em}
    \textbf{Coming up on Wednesday:}
    \begin{itemize}
        \item \textbf{Plug-and-Play (PnP) Priors \& RED:} How to use pre-trained, off-the-shelf denoisers as universal regularisers, entirely decoupling the neural network from the forward physics.
    \end{itemize}
\end{frame}

\begin{frame}{Preview: Wednesday and Thursday}
    What if we want infinite depth without running out of GPU memory? How do we handle uncertainty?

    \vspace{0.5em}
    \textbf{Coming up on Thursday:}
    \begin{itemize}
        \item \textbf{Deep Equilibrium Models (DEQs):} Pushing unrolled networks to infinite depth $K \to \infty$ using root-finding algorithms and revisiting implicit differentiation.
        \item \textbf{Generative Models:} Autoencoders, VAEs, and Score-Based Diffusion Models for sampling from the posterior distribution.
    \end{itemize}
\end{frame}

\begin{frame}[allowframebreaks]{References}
    \bibliographystyle{plainnat}
    \bibliography{references}
\end{frame}

\end{document}